\documentclass[12pt,a4paper]{article}
\usepackage{fontspec}
\usepackage{polyglossia}
\setmainfont{TeX Gyre Termes}
\newfontfamily\cyrillicfont{TeX Gyre Termes}
\setdefaultlanguage{ukrainian}
\setotherlanguage{english}
\usepackage{amsmath,amssymb,geometry,graphicx}
\geometry{margin=25mm}
\setlength{\parindent}{1.25cm}
\setlength{\parskip}{0.25em}
\begin{document}
\title{2022 АНАЛІЗ ПОВЕДІНКИ ПОХИБКИ ОБЧИСЛЕНЬ ... ФУНКЦІЇ РУНГЕ (ВАК)}
\maketitle
2022 ￿￿￿￿￿￿ ￿￿￿￿￿￿￿￿￿ ￿￿￿￿￿￿￿ ￿￿￿￿￿￿￿￿￿ ... ￿￿￿￿￿￿￿ ￿￿￿￿￿ (￿￿￿) 27 ￿￿￿￿￿￿ 2026 ￿. 2022 ￿￿￿￿￿￿ ￿￿￿￿￿￿￿￿￿ ￿￿￿￿￿￿￿ ￿￿￿￿￿￿￿￿￿ ... ￿￿￿￿￿￿￿ ￿￿￿￿￿ (￿￿￿) 27 ￿￿￿￿￿￿ 2026 ￿. 159 ￿￿￿ 514.18 ￿￿￿￿￿￿ ￿￿￿￿￿￿￿￿￿ ￿￿￿￿￿￿￿ ￿￿￿￿￿￿￿￿￿ ￿￿￿ ￿￿￿￿-￿￿￿￿￿￿￿￿￿￿￿￿ ￿￿￿￿￿￿￿ ￿￿￿￿￿ ￿￿￿￿￿￿￿￿￿ ￿.￿., ￿.￿.￿., sulico@gmail.com, ORCID: 0000-0002- 1953-0410 ￿￿￿￿￿￿￿￿￿ ￿.￿., ￿.￿.￿., o.zalevska@kpi.ua, ORCID: 0000-0002-3163-1695 ￿￿￿￿￿￿￿￿￿￿￿ ￿.￿. o.zalevska@kpi.ua, ORCID: 0000-0003-4673-3894 ￿￿￿￿￿￿￿￿￿￿￿￿ ￿￿￿￿￿￿￿￿￿ ￿￿￿￿￿￿￿￿￿￿￿ ￿￿￿￿￿￿￿ «￿￿￿￿￿￿￿￿￿ ￿￿￿￿￿￿￿￿￿￿￿￿￿ ￿￿￿￿￿￿￿￿ ￿￿￿￿￿ ￿￿￿￿￿ ￿￿￿￿￿￿￿￿￿￿￿» (￿￿￿￿￿￿￿) ￿￿￿￿￿￿￿￿￿ ￿.￿., ￿.￿.￿., spirintsev@gmail.com, ORCID: 0000-0001-5728-6626 ￿￿￿￿￿￿￿￿￿￿￿￿￿￿￿ ￿￿￿￿￿￿￿￿￿ ￿￿￿￿￿￿￿￿￿￿￿￿ ￿￿￿￿￿￿￿￿￿￿￿ ￿￿￿￿￿ ￿￿￿￿￿￿￿ ￿￿￿￿￿￿￿￿￿￿￿￿￿ (￿￿￿￿￿￿￿) ￿ ￿￿￿￿￿￿ ￿￿￿￿￿￿￿￿￿￿￿￿￿ ￿￿￿￿￿￿￿ ￿￿￿￿￿￿￿￿￿ ￿￿￿ ￿￿￿￿￿￿￿￿￿￿￿￿ ￿￿ ￿￿￿￿￿￿￿ ￿￿￿￿￿ ￿￿￿￿ -￿￿￿￿￿￿￿￿￿￿￿￿. ￿￿￿￿￿￿￿￿￿ ￿￿￿￿￿￿￿ ￿￿￿￿￿￿￿￿￿￿￿￿ ￿￿ ￿￿￿￿￿￿￿ ￿￿￿￿￿￿￿￿￿￿￿￿￿ ￿￿￿￿￿￿￿ ￿￿￿￿￿￿￿￿￿￿￿￿ ￿￿￿￿￿￿￿￿￿￿￿. ￿￿￿￿￿￿￿￿￿ ￿ ￿￿￿￿￿￿￿￿ ￿￿￿￿￿￿ ￿￿￿￿’￿￿￿￿￿￿￿ ￿￿￿￿￿￿￿￿ ￿￿￿￿￿￿￿￿￿￿￿￿ ￿￿￿￿￿￿ ￿￿￿￿￿￿ ￿￿￿￿- ￿￿￿￿￿￿￿ ￿￿ ￿￿￿￿￿￿￿￿￿￿￿￿ ￿￿￿￿￿￿￿￿￿￿ ￿ ￿￿￿￿￿￿￿￿￿ ￿￿￿￿￿￿ ￿￿￿￿￿￿￿￿. ￿ ￿￿￿ ￿￿￿ ￿￿￿￿￿￿￿￿￿ ￿￿￿￿￿￿￿￿￿￿￿￿ ￿￿￿￿￿￿￿￿￿￿￿￿ ￿￿￿￿￿ ￿￿￿￿￿￿￿ ￿￿￿￿￿￿, ￿￿￿ ￿￿￿￿￿￿￿￿￿￿￿￿ ￿￿￿￿￿￿￿￿￿￿ ￿￿￿￿￿￿￿ ￿￿ ￿￿￿￿￿￿￿￿￿￿￿ , ￿￿￿￿￿￿￿￿￿￿ ￿￿￿￿￿, ￿￿￿￿￿￿￿￿￿￿￿, ￿￿￿￿￿￿￿￿￿ ￿￿￿￿￿￿￿￿, ￿￿ ￿￿￿￿ ￿ ￿￿￿￿￿ ￿￿￿￿￿￿￿ ￿￿￿￿￿ ￿ ￿￿ ￿￿￿￿￿￿￿￿￿ ￿￿￿￿￿￿￿ ￿ ￿￿￿￿￿￿￿￿￿. ￿￿￿￿￿￿ ￿￿￿￿￿￿￿￿￿￿￿￿￿ ￿ ￿￿￿￿￿￿￿￿￿￿￿￿￿ ￿￿￿￿￿￿￿￿￿￿￿￿ ￿ ￿￿￿￿￿￿￿￿￿￿￿￿￿￿￿￿￿ ￿￿￿￿￿￿￿￿ ￿￿￿￿￿ ￿￿￿￿￿￿￿￿￿￿￿￿￿ ￿￿￿￿￿.
\par
\[￿￿￿￿￿￿￿￿￿￿￿￿￿￿￿￿￿￿￿￿￿￿￿￿￿￿￿￿￿￿￿￿￿￿￿.￿￿￿￿￿￿￿￿￿￿￿￿￿￿￿￿￿￿￿￿￿￿￿￿￿￿￿￿ - ￿￿￿￿￿￿￿￿￿￿￿￿￿￿￿￿￿￿￿￿￿￿￿￿￿￿￿￿￿￿￿￿￿￿￿￿￿￿￿￿, ￿￿￿￿￿￿￿￿￿, 460\times\]
800 ￿￿￿￿￿￿￿￿ ￿ ￿￿￿￿￿ ￿￿￿￿￿￿￿, ￿￿ ￿￿￿￿￿￿ ￿￿￿￿￿￿￿￿ ￿￿￿￿￿￿￿￿￿￿￿￿ ￿￿￿￿￿￿￿￿￿￿ ￿ ￿￿￿￿￿￿￿￿￿￿￿ ￿￿￿￿￿ ￿￿￿￿￿￿￿ ￿￿￿￿￿￿￿￿￿￿ ￿￿￿￿￿￿￿￿￿. ￿￿ ￿￿￿￿￿￿￿, ￿￿ ￿￿￿￿￿￿￿￿￿￿￿￿ ￿￿￿￿￿￿￿￿￿￿, ￿￿ ￿￿￿ ￿￿￿￿￿￿￿￿￿￿, ￿￿￿ ￿ ￿￿￿ ￿￿￿￿￿￿￿￿￿ ￿￿￿￿￿￿￿￿￿￿, ￿￿￿￿￿￿￿￿￿, ￿￿￿￿￿￿￿￿￿￿￿ ￿￿ ￿￿￿￿￿￿￿￿￿ ￿￿￿￿￿￿￿￿￿￿￿￿. ￿￿￿￿￿￿￿￿￿￿￿￿ ￿￿￿￿￿￿ ￿￿￿￿￿ ￿ ￿ ￿￿￿￿￿￿￿￿ ￿￿￿￿￿￿￿￿. ￿￿￿￿￿￿￿￿ ￿￿￿ ￿￿￿￿￿￿￿￿￿￿ ￿￿￿￿￿￿￿￿￿￿￿￿ ￿￿￿￿￿￿￿￿ ￿￿￿￿￿￿￿￿￿￿ ￿￿￿￿￿￿￿￿ ￿￿￿￿￿￿￿￿￿￿￿￿ ￿￿￿￿￿￿￿ ￿￿￿￿￿￿￿￿. ￿￿￿ ￿￿￿￿￿￿￿ ￿￿￿￿￿￿￿￿￿￿￿￿ ￿￿￿￿￿￿ ￿￿￿￿￿￿￿￿ ￿￿￿￿ ￿￿￿￿￿￿￿￿ ￿￿￿￿￿ ￿￿￿￿￿￿￿￿. ￿￿￿, ￿￿￿ ￿￿￿￿￿ ￿￿￿￿￿￿￿￿￿￿￿￿￿￿￿￿ ￿￿￿ ￿￿￿￿￿￿￿￿￿ ￿￿￿￿￿￿ ￿￿￿￿￿￿￿￿ ￿￿￿￿￿￿￿ ￿￿ ￿￿ ￿￿￿￿￿￿￿￿￿￿￿ ￿￿￿￿￿￿￿￿￿￿, ￿￿￿ ￿￿￿ ￿￿￿￿￿￿￿ ￿￿￿￿￿￿￿￿￿￿￿￿￿￿￿￿￿ ￿￿￿￿￿ ￿￿ ￿￿￿￿￿￿ ￿￿￿￿￿￿￿￿ ￿￿￿￿ ￿￿￿￿￿￿￿￿￿￿￿￿￿￿￿. ￿￿￿￿￿ ￿￿￿￿￿￿￿￿￿￿ ￿ ￿￿￿￿￿￿ ￿￿￿￿￿￿￿￿￿ ￿￿￿￿￿￿￿￿￿￿￿￿ ￿￿￿￿￿￿ ￿￿￿￿￿￿￿￿ ￿￿ ￿￿￿￿￿￿￿￿ ￿￿￿￿￿￿￿ ￿￿￿￿￿. ￿ ￿￿￿￿￿￿ ￿￿￿￿￿￿￿￿￿￿￿￿￿ ￿￿￿￿￿￿￿￿￿￿￿￿￿ ￿￿ 160 ￿￿￿￿￿￿￿ ￿￿ ￿￿￿￿￿￿￿￿￿ ￿￿￿￿￿ ￿￿￿￿￿ ￿￿￿￿￿￿￿ ￿￿￿￿￿, ￿￿￿￿￿￿￿￿ ￿￿￿￿￿￿￿￿ ￿￿￿￿￿ ￿￿￿￿￿￿￿ ￿￿￿ ￿￿￿￿’￿￿￿￿￿￿ ￿￿￿￿￿￿￿￿￿ ￿￿￿￿￿￿ ￿￿￿￿￿- ￿￿￿￿￿￿￿ ￿￿￿￿￿￿￿ ￿￿￿￿￿. ￿￿￿￿￿￿￿ ￿￿￿￿￿: ￿￿￿￿￿￿￿￿￿￿￿￿, ￿￿￿￿￿￿￿￿￿￿￿￿￿￿ ￿￿￿￿￿￿￿ ￿￿￿￿￿, ￿￿￿￿￿￿￿ ￿￿￿￿￿￿￿￿￿￿￿￿, ￿￿￿￿￿￿￿ ￿￿￿￿￿. ￿￿￿￿￿￿￿￿￿￿ ￿￿￿￿￿￿￿￿. ￿￿ ￿￿￿￿￿￿￿￿￿ ￿￿￿￿￿￿￿ ￿￿￿￿￿￿￿￿￿￿￿￿ ￿￿￿￿￿ ￿￿￿￿￿￿￿￿ ￿￿￿￿￿ ￿￿￿￿￿￿￿￿. ￿￿￿ ￿￿ ￿￿￿￿￿￿ ￿￿￿￿￿￿￿￿￿ ￿￿￿￿￿￿￿￿￿ ￿￿￿￿ ￿￿￿￿￿￿￿￿￿￿￿￿ ￿￿ -￿￿ ￿￿￿￿￿￿￿ ￿￿￿￿￿￿￿ ￿￿￿ ￿￿￿￿￿￿￿￿ ￿￿￿￿￿￿ ￿￿￿￿￿￿￿ ￿￿￿￿￿. ￿￿￿￿￿ ￿￿￿￿￿￿ ￿￿￿￿￿￿￿ ￿￿￿ 0 ￿ ￿￿￿￿￿￿￿ ￿￿￿￿￿￿ (￿￿￿￿￿￿￿￿￿, ￿￿￿ ￿￿￿￿’￿￿￿￿￿￿ ￿￿￿￿￿￿￿￿￿￿￿￿￿￿ ￿￿￿￿￿￿￿￿ ￿￿￿￿￿￿ ￿￿￿ ￿￿￿￿￿￿￿￿￿ ￿￿￿￿￿￿￿￿￿￿￿￿￿ ￿￿￿￿￿￿￿. ￿￿￿ ￿￿￿￿￿￿￿, ￿￿ ￿￿￿￿￿￿￿￿￿￿￿￿￿￿, ￿￿￿￿￿￿￿ ￿￿￿￿￿￿ ￿￿￿￿￿￿￿ ￿, ￿￿￿￿￿￿￿￿￿, ￿￿￿￿￿￿￿￿￿ ￿￿￿￿￿￿￿￿￿￿￿￿ ￿￿ COVID -19 ￿￿ ￿￿￿￿￿￿ ￿￿￿￿￿￿ ￿￿￿￿. ￿￿￿￿￿￿￿ ￿￿ ￿￿￿￿￿￿￿ ￿￿￿￿￿￿ ￿￿￿￿￿￿￿￿ ￿￿￿￿￿￿ ￿￿￿￿￿￿￿￿ ￿￿￿￿￿￿￿￿￿, ￿￿￿ ￿￿￿￿￿￿￿￿￿ ￿￿￿￿￿￿ ￿￿￿￿￿￿￿￿ ￿ ￿￿￿￿￿￿￿￿￿￿￿ ￿￿￿￿￿￿￿ ￿￿￿￿￿￿￿ ￿￿￿￿￿ ￿ ￿￿￿￿￿￿￿ ￿￿￿￿￿￿￿￿￿￿￿￿ ￿￿￿ ￿￿￿￿￿￿￿ ￿￿￿￿￿￿￿. ￿￿￿￿￿￿￿￿￿ ￿￿ ￿￿￿￿￿￿￿ ￿￿￿￿￿￿￿￿￿ ￿￿￿￿￿￿￿ ￿￿￿ ￿￿￿￿￿￿￿￿￿￿ ￿￿￿￿￿￿ ￿￿￿￿￿￿￿￿￿￿￿￿ ￿ ￿￿￿￿￿￿￿ ￿￿￿￿￿. ￿￿￿ ￿￿￿￿￿￿￿￿￿￿￿ ￿￿￿￿￿￿￿￿￿ ￿￿￿￿￿￿￿ ￿￿￿￿ ￿￿￿￿￿￿￿￿￿￿￿￿￿ ￿￿￿ ￿￿￿￿’￿￿￿￿￿￿ ￿￿￿￿ ￿￿￿￿￿￿ ￿￿￿￿￿￿￿￿￿￿￿ ￿￿￿￿￿ - ￿￿￿￿￿￿￿￿￿￿￿￿. ￿￿￿￿￿￿ ￿￿￿￿￿￿￿￿ ￿￿￿￿￿￿￿￿￿￿ ￿ ￿￿￿￿￿￿￿￿￿￿ . ￿ ￿￿￿￿￿￿ [1] ￿￿￿￿￿￿￿￿ ￿￿￿￿ ￿￿￿￿￿￿￿￿ ￿￿￿￿￿, ￿￿￿￿￿￿￿￿￿￿￿￿￿ ￿￿￿￿￿ ￿￿￿￿￿￿￿￿￿￿ ￿￿￿￿￿￿￿ ￿￿￿￿￿￿￿ ￿￿￿ ￿￿￿￿￿￿￿￿￿￿￿￿ ￿￿￿￿￿￿￿ ￿￿￿￿￿￿￿￿. ￿ ￿￿￿￿￿￿ [2] ￿￿￿￿￿￿ ￿￿￿￿￿￿￿ ￿￿￿￿￿￿￿￿￿￿￿ ￿￿￿￿￿￿￿￿ ￿￿￿￿￿￿￿ ￿￿￿￿￿￿￿ ￿￿ ￿￿￿￿￿￿￿￿￿ ￿￿￿￿￿￿￿￿￿￿￿￿ ￿￿￿￿￿￿￿￿￿ ￿￿￿￿￿￿￿￿. ￿￿￿￿￿ ￿￿￿￿￿￿￿￿￿￿￿ ￿￿￿￿￿￿￿￿￿￿￿￿ ￿￿￿￿￿ ￿￿￿￿￿￿￿￿￿￿￿, ￿ ￿￿￿￿￿￿ ￿￿￿￿￿ ￿￿ ￿￿￿￿￿￿￿￿￿￿￿￿￿￿￿￿ ￿￿￿￿￿￿￿￿, ￿￿ ￿￿￿￿￿￿￿￿ ￿￿￿￿￿￿￿￿￿￿ ￿￿￿￿￿￿￿￿￿￿￿ ￿￿￿￿￿￿ ￿￿￿￿￿￿￿￿￿￿￿￿ ￿￿￿ ￿￿￿￿￿ ￿￿￿￿￿￿￿￿ ￿￿￿￿￿ ￿￿￿￿￿ ￿￿￿￿￿￿￿￿￿￿ ￿ ￿￿￿￿￿￿ [3]. ￿ [4] ￿￿￿￿￿￿￿￿￿￿￿￿￿ ￿￿￿￿￿￿￿￿￿￿￿￿ ￿￿￿￿￿￿￿￿￿￿ ￿￿￿￿￿ ￿￿￿￿￿￿￿￿￿￿￿￿ ￿￿￿￿￿￿￿￿￿￿ ￿￿￿￿￿￿￿￿￿￿, ￿￿ ￿￿￿￿￿ ￿￿￿ ￿￿￿￿￿￿￿￿￿￿ ￿￿￿￿￿￿ ￿￿￿￿￿￿￿￿ ￿￿￿￿￿￿￿. ￿￿￿￿￿￿￿ ￿￿￿￿ -￿￿￿￿￿￿￿￿￿￿￿￿ ￿￿￿￿￿￿ ￿ ￿￿￿￿￿￿ [5], ￿￿￿￿￿￿￿￿￿ ￿￿￿￿￿￿￿ ￿￿￿￿￿￿ ￿￿￿￿￿￿￿￿￿￿￿￿￿￿￿ ￿￿￿￿￿￿￿ ￿￿￿￿￿ ￿￿ ￿￿ ￿￿￿￿￿￿￿￿￿, ￿ ￿￿￿￿ ￿￿￿￿￿￿￿￿￿￿￿￿￿ ￿￿ ￿￿￿￿￿￿￿￿. ￿ ￿￿￿￿￿￿ [6] ￿￿￿￿￿￿￿￿￿ ￿￿￿￿￿￿ ￿￿￿￿￿￿ ￿￿￿￿￿￿￿￿￿￿￿￿￿￿￿ ￿￿￿￿￿￿￿ ￿￿￿￿￿ ￿￿ ￿￿￿￿￿￿￿￿￿￿￿￿ ￿￿￿￿￿￿￿￿￿￿￿ ￿￿￿￿￿￿￿￿ , ￿￿￿￿￿￿￿￿￿￿￿ ￿￿￿￿￿￿￿ ￿￿￿￿￿￿￿￿￿. ￿￿￿￿￿￿￿￿ ￿￿￿￿￿￿￿￿￿ ￿￿￿￿￿￿￿ ￿￿￿￿￿￿￿￿￿￿￿￿ ￿￿￿￿￿￿￿￿ ￿￿￿￿￿ ￿￿ ￿￿￿￿￿ ￿￿￿￿￿￿￿￿￿ ￿￿ ￿￿￿￿￿￿￿￿￿ ￿￿￿￿￿￿￿￿￿ ￿￿￿￿￿￿￿￿￿￿ ￿￿￿￿￿￿￿￿, ￿￿￿ ￿￿￿￿￿￿ ￿￿￿￿￿￿￿￿￿￿￿ ￿￿￿￿￿￿￿￿￿￿￿￿ ￿￿￿￿￿ ￿￿￿￿￿￿￿￿￿, ￿￿￿￿￿￿￿￿ ￿ ￿￿￿￿￿￿ [7]. ￿￿￿￿￿￿￿￿￿￿￿￿ ￿￿￿￿￿ ￿￿￿￿￿￿. ￿￿￿￿￿ ￿￿￿￿￿￿￿￿￿￿￿ ￿ ￿￿￿￿￿￿ ￿￿￿￿￿￿￿￿￿￿￿ ￿￿￿￿￿￿￿￿￿￿￿￿ ￿￿￿￿￿￿￿ ￿￿￿￿￿ ￿￿ 1 1
\par
￿￿￿￿￿￿￿￿￿ ￿￿￿￿￿ ￿￿￿￿￿ ￿￿￿￿￿￿￿￿￿￿￿￿￿￿￿ ￿￿￿￿￿￿￿ ￿￿￿￿￿ ￿￿ ￿￿￿￿￿￿￿￿￿￿ ￿￿￿￿￿￿￿￿￿￿￿ ￿ ￿￿￿￿￿￿￿ ￿￿￿￿￿￿￿￿. ￿￿￿￿￿￿￿ ￿￿￿￿￿￿￿. ￿ ￿￿￿￿￿￿￿￿￿￿ ￿￿￿ ￿￿￿￿￿￿ ￿￿￿￿￿￿￿ ￿￿￿￿￿ ￿￿￿￿￿￿￿ ￿￿￿￿￿￿￿￿ ￿￿￿￿￿￿￿￿, ￿￿￿ ￿￿ ￿￿￿￿￿￿￿￿￿ ￿￿￿￿￿￿￿￿￿￿￿￿ ￿￿￿￿￿￿￿￿￿￿￿￿ ￿￿￿￿￿ ￿￿￿￿￿￿￿ ￿￿￿￿￿￿￿￿￿ ￿￿￿￿￿￿￿￿￿. ￿￿￿ ￿￿￿￿’￿￿￿￿￿￿ ￿￿￿￿ ￿￿￿￿￿￿￿￿￿ ￿￿￿￿￿ , ￿￿￿￿, ￿￿￿￿￿￿￿￿￿, ￿￿￿￿ ￿￿￿￿￿￿￿￿ ￿ ￿￿￿￿￿￿￿￿￿, ￿ ￿￿￿￿￿ ￿￿￿￿￿￿￿￿￿￿￿ ￿￿￿￿￿￿￿ 161 ￿￿￿￿￿￿￿ ￿￿￿￿￿￿￿ ￿￿￿ ￿￿￿￿￿￿￿￿￿￿￿￿, ￿ ￿￿￿￿￿￿, ￿￿￿￿￿￿￿￿￿￿￿ ￿￿￿￿￿￿ ￿￿￿￿ ￿￿￿￿￿￿￿￿, ￿ ￿￿￿￿￿￿￿￿￿ ￿￿￿￿￿￿￿￿￿￿, ￿￿￿￿￿￿￿ ￿￿￿￿￿￿￿￿￿￿ ￿￿￿￿￿￿￿ ￿￿ ￿￿￿￿ ￿￿￿￿￿￿, ￿￿￿￿￿ ￿￿￿￿￿￿￿￿ ￿￿￿￿￿￿￿￿￿ ￿￿￿￿￿￿￿￿￿￿ ￿￿￿￿’￿￿￿￿￿￿ ￿ ￿￿￿￿￿￿￿￿￿￿ ￿￿￿￿￿￿￿￿, ￿￿ ￿￿￿￿’￿￿￿￿￿￿ ￿￿￿￿￿ ￿ ￿￿￿￿￿￿￿￿ ￿￿￿￿￿￿￿￿￿ ￿￿￿￿￿￿￿￿ ￿￿￿￿￿, ￿￿￿￿￿ ￿￿ ￿￿￿￿￿￿￿￿￿ ￿￿￿￿￿￿￿￿￿, ￿￿￿ ￿￿￿￿￿￿￿￿￿￿ ￿￿￿￿￿￿￿￿￿. ￿￿￿￿￿￿ ￿￿￿￿’￿￿￿￿￿￿ ￿￿￿￿￿￿ ￿￿￿￿￿￿￿￿￿￿￿￿ ￿￿￿￿￿ ￿￿￿￿￿￿￿ ￿￿￿￿, ￿￿￿ ￿￿￿￿￿￿￿￿ ￿￿￿￿￿￿ ￿￿￿￿ ￿￿￿￿￿￿ ￿￿￿￿￿￿￿￿￿￿￿￿￿￿￿￿ ￿￿ ￿￿￿￿￿￿￿￿. ￿￿￿￿￿ ￿ ￿￿￿￿￿￿￿￿￿ ￿￿￿￿￿￿￿ ￿ ￿￿￿￿￿ ￿￿￿￿￿￿￿￿. ￿￿￿￿￿￿￿ ￿￿￿￿￿￿￿￿ ￿￿￿ ￿￿￿￿￿ ￿￿￿￿￿￿: ￿ ￿ ￿ ￿ ￿ ￿￿ ￿ ￿ ￿ ￿ ￿￿ ￿ ￿ ￿￿ ￿ ￿ ￿ 1 2 1 1 1 1 2 1 1 ()n i i n i ii ni i i i i i n x x x x x x x x x xPx x x x x x x x x x yx ￿￿ ￿ ￿￿ ￿ ￿ ￿ ￿ ￿ ￿ ￿￿ ￿ ￿ ￿ ￿ ￿ ￿ ￿￿ . ￿￿￿￿￿ ￿￿￿￿￿￿￿￿ ￿￿￿ ￿￿￿￿ ￿￿￿￿￿￿￿￿ ￿ ￿￿￿￿￿￿￿￿. ￿￿￿￿￿￿￿￿ ￿￿￿￿￿￿￿￿￿ ￿￿￿￿￿ ￿￿￿￿￿￿ ￿ ￿￿, ￿￿ ￿￿￿￿￿￿￿ ￿￿￿￿￿￿￿￿￿ ￿￿￿￿￿ ￿￿￿￿￿￿￿ ￿￿￿ ￿￿￿￿￿￿￿￿￿￿ ￿￿￿￿￿￿￿￿￿ ￿￿￿￿￿￿ ￿￿￿￿￿￿￿￿￿￿￿￿. ￿￿￿ ￿￿￿￿￿￿￿ ￿￿￿￿￿￿￿ ￿ ￿￿￿ ￿￿￿￿￿￿￿ ￿￿￿￿￿￿￿￿ ￿￿￿￿￿. ￿￿ ￿￿￿￿￿￿￿￿ ￿￿￿￿￿ ￿￿￿￿￿￿ ￿￿￿￿￿￿￿ ￿￿￿￿￿: ￿￿￿￿￿ ￿￿￿￿￿￿, ￿￿ ￿￿￿ ￿￿￿￿￿￿ ￿￿￿￿ ￿￿￿￿￿￿ ￿￿￿￿￿￿ ￿￿￿￿￿￿￿ ￿￿￿￿￿￿￿ ￿￿￿￿￿￿￿￿￿￿￿￿ ￿￿￿￿￿￿ ￿￿ ￿￿￿￿￿￿￿￿￿￿￿￿￿￿, ￿￿￿￿ ￿￿￿￿￿￿￿￿￿ ￿￿￿￿￿￿ , ￿ ￿ ￿￿￿ ￿ ￿￿￿￿￿￿￿ ￿￿￿￿￿￿￿￿, ￿￿￿￿￿￿￿, ￿￿ ￿￿￿￿￿￿￿￿ ￿￿ ￿￿￿. 1. ￿￿￿. 1. ￿￿￿￿￿￿￿ ￿￿￿￿￿ 162 ￿ ￿￿￿￿￿ ￿￿￿￿￿￿ ￿￿￿￿ ￿￿￿￿￿￿￿￿￿￿ ￿￿￿ ￿￿￿ ￿￿￿￿￿￿￿￿￿￿￿￿ ￿￿￿￿￿￿￿ ￿￿￿￿￿ ￿￿￿￿￿￿￿￿￿￿￿ ￿￿￿￿-￿￿￿￿￿￿￿ ￿￿￿￿￿￿ ￿￿￿￿￿￿￿: 222 12 ()( ) ( ) 12( ) . nxxx x x x nG x y e y e y e ￿￿￿ ￿￿￿ ￿ ￿ ￿￿ ￿ ￿ ￿ ￿￿￿ ￿￿￿￿￿￿￿￿￿ ￿￿￿￿￿￿￿￿￿￿￿￿￿￿ ￿￿￿￿￿￿ ￿￿￿￿￿￿￿ ￿￿￿￿￿ ￿￿￿￿￿ ￿￿￿￿￿￿￿￿ ￿￿￿￿￿￿￿ ￿￿￿￿￿￿ ￿￿￿￿￿￿￿, ￿￿ ￿￿￿￿￿ ￿ ￿￿￿￿￿￿￿ ￿￿￿￿ ￿￿￿￿￿￿￿￿￿￿￿￿￿￿￿ ￿￿￿￿￿￿￿￿ ￿￿￿￿￿: ￿ ￿ ￿ ￿ ￿ ￿ ￿ ￿ x , y ￿￿￿￿ ￿￿￿ i i ti iY t X ￿￿￿￿￿￿￿￿￿￿￿ ￿￿￿￿￿￿￿￿ ￿ ￿￿ ￿￿￿￿￿￿￿ ￿￿￿￿￿￿￿￿￿￿ ￿￿ ￿￿￿￿￿￿￿￿: 2 max min ( 1) ,() ￿￿ ￿￿ ￿ n xx ￿￿ ￿max, xmin — ￿￿￿￿￿￿￿￿￿￿￿ ￿￿ ￿￿￿￿￿￿￿￿￿￿ ￿￿￿￿￿￿￿￿ ￿￿￿￿￿￿￿￿￿ ￿. ￿￿￿￿￿￿￿￿￿￿￿￿ ￿￿￿￿￿￿￿￿￿￿￿ ￿￿￿￿￿￿￿￿￿￿￿￿ ￿￿ ￿￿￿￿￿￿￿￿￿￿ ￿￿￿￿￿￿￿ ￿￿￿￿￿￿￿￿￿ ￿￿￿￿￿￿￿￿￿￿￿ ￿￿ ￿￿￿￿￿￿￿￿￿ ￿￿￿￿’￿￿￿￿￿￿￿ ￿￿￿￿￿￿￿￿, ￿￿￿￿￿￿￿￿￿ ￿ ￿￿￿￿￿￿￿￿￿￿￿ ￿￿￿￿￿￿￿￿￿￿ Microsoft Visual Studio 2020, ￿￿￿￿￿ ￿\#. ￿￿￿. 2. ￿￿￿￿￿￿￿￿￿￿￿￿ ￿￿￿￿￿￿￿ ￿￿￿￿￿ ￿￿￿￿￿￿￿￿￿ ￿￿￿￿￿￿￿￿ ￿￿ ￿￿￿ . 2 ￿￿￿￿￿￿￿￿￿￿￿￿ ￿￿￿￿￿￿￿￿￿￿ ￿￿￿￿￿￿￿￿￿￿￿￿ ￿￿￿￿￿￿￿ ￿￿￿￿￿ ￿￿ ￿￿￿￿￿￿￿ ￿￿￿￿￿￿￿￿￿￿￿￿ ￿￿￿ 20 ￿￿￿￿￿￿ ￿￿￿￿￿￿￿ ￿￿￿￿￿￿￿ ￿￿￿￿￿￿￿￿. ￿￿￿￿￿￿￿￿, ￿￿ ￿￿￿￿￿￿￿￿￿￿ ￿￿ ￿￿￿￿￿￿ ￿￿￿￿ ￿￿￿￿￿￿￿￿￿ ￿￿￿ ￿￿￿￿￿￿￿￿￿ ￿￿￿￿￿￿￿￿￿￿￿, ￿￿￿￿￿￿￿￿ ￿￿ ￿￿￿￿￿￿ ￿￿￿￿￿￿￿￿, ￿￿ ￿￿￿￿￿￿￿￿￿￿￿ ￿￿￿￿￿￿￿ ￿￿￿￿￿. 163 ￿￿ ￿￿￿ . 3 ￿￿￿￿￿￿￿￿ ￿￿￿￿￿￿￿￿￿ ￿￿￿￿￿￿￿￿￿￿￿￿ ￿￿ ￿￿￿￿￿￿￿￿￿ ￿￿￿￿￿￿￿￿￿ ￿￿￿￿￿￿￿ ￿￿￿￿￿. ￿￿￿￿￿￿￿￿￿ ￿￿￿￿￿￿￿￿￿￿￿￿ ￿￿￿￿￿￿ ￿￿￿￿￿￿ ￿￿ ￿￿￿￿￿ ￿￿￿￿￿￿￿￿, ￿￿￿ ￿￿ ￿￿￿￿￿￿ ￿￿￿￿￿￿￿￿ ￿￿￿ ￿￿ ￿￿￿￿￿￿￿￿ ￿￿￿￿￿￿ ￿￿￿￿￿￿￿￿￿￿. ￿￿￿. 3 ￿￿￿￿￿￿￿￿￿￿￿￿ ￿￿￿￿￿￿￿ ￿￿￿￿￿ ￿￿￿￿￿￿￿￿￿ ￿￿￿￿-￿￿￿￿￿￿￿￿ ￿￿ ￿￿￿ . 4 ￿￿￿￿￿￿￿￿ ￿￿￿￿￿￿￿￿￿￿ ￿￿￿￿￿￿￿￿￿￿￿￿ ￿￿ ￿￿￿￿￿￿￿￿￿ ￿￿￿￿￿￿￿￿￿￿￿￿￿ ￿￿￿￿￿￿￿ ￿￿￿￿￿. ￿￿￿. 4. ￿￿￿￿￿￿￿￿￿￿￿￿ ￿￿￿￿￿￿￿ ￿￿￿￿￿ ￿￿￿￿￿￿￿￿￿￿￿￿￿ ￿￿￿￿-￿￿￿￿￿￿￿￿ 164 ￿￿￿￿￿￿￿ ￿￿￿￿￿￿￿￿￿￿￿￿ ￿￿ ￿￿￿￿￿￿￿￿ ￿￿￿￿￿￿￿￿ ￿￿￿￿￿, ￿￿￿ ￿￿￿ ￿￿￿￿￿￿￿￿￿￿￿￿ ￿￿￿￿￿￿￿ ￿￿￿￿￿￿￿￿. ￿￿ ￿￿￿￿￿￿ ￿￿￿￿￿￿￿￿ ￿￿￿￿￿ ￿￿￿￿￿￿￿ ￿￿￿￿￿￿. ￿￿￿￿￿￿￿￿￿￿ ￿￿￿￿￿￿￿￿￿￿￿￿ ￿￿ ￿￿￿￿￿￿￿￿￿ ￿￿￿￿￿￿￿￿ ￿￿￿￿￿￿￿ ￿￿￿￿￿ ￿￿￿￿￿￿￿￿ ￿￿ ￿￿￿. 5. ￿￿￿. 5. ￿￿￿￿￿￿￿￿￿￿￿￿ ￿￿￿￿￿￿￿ ￿￿￿￿￿ ￿￿￿￿￿￿￿￿ ￿￿￿￿-￿￿￿￿￿￿￿￿ ￿￿ ￿￿￿￿￿ ￿ ￿ ￿￿￿￿￿￿￿, ￿ ￿ ￿￿￿￿￿￿￿￿￿, ￿￿￿ ￿￿￿￿￿￿￿￿￿￿￿￿ ￿￿ ￿￿￿￿￿￿, ￿￿￿￿￿￿￿ ￿￿￿￿￿￿￿￿￿￿￿￿ ￿ ￿￿￿￿￿￿￿￿￿￿￿. ￿￿￿￿￿ ￿￿￿￿￿, ￿ ￿￿￿￿￿￿ ￿￿￿￿￿ ￿￿￿￿￿￿￿ ￿￿￿￿￿￿￿￿ ￿￿￿￿￿￿￿￿ ￿￿￿￿￿. ￿￿￿￿￿￿￿ ￿￿￿￿￿￿￿￿￿ ￿￿￿￿￿ ￿￿ ￿￿￿￿￿￿￿￿. ￿ ￿￿￿￿’￿￿￿￿￿￿￿ ￿￿￿￿￿￿￿ ￿￿￿￿￿￿￿￿￿￿￿ ￿￿￿￿￿￿￿￿￿￿ ￿￿￿￿￿￿￿￿￿ ￿￿￿￿￿￿￿￿￿￿ α. ￿ ￿￿￿￿￿￿￿ ￿￿￿￿￿￿￿￿ ￿￿￿￿ ￿￿￿￿￿￿ ￿ ￿￿￿￿￿￿￿￿￿ ￿￿￿ ￿￿￿￿￿ ￿￿￿￿￿ ￿￿￿￿￿￿￿￿￿. ￿￿￿￿￿ ￿￿￿￿￿ ￿￿￿￿￿￿ ￿￿￿￿￿￿￿￿ ￿￿￿￿￿￿￿￿￿￿￿￿￿ ￿￿￿￿￿￿￿ α ￿ ￿￿￿￿￿￿￿￿￿ ￿￿￿￿￿￿￿￿￿￿￿￿ ￿￿￿￿￿￿￿￿. ￿￿￿￿￿￿￿￿. ￿￿￿￿ ￿￿￿￿￿￿￿￿￿ ￿￿￿￿￿￿ ￿￿￿￿￿￿￿￿￿ ￿￿￿￿￿￿￿￿￿￿￿￿ ￿￿￿￿￿￿ ￿￿￿￿￿￿￿￿ ￿￿ ￿￿￿￿￿￿￿￿ ￿￿￿￿￿￿￿ ￿￿￿￿￿ , ￿￿ ￿￿￿ ￿￿￿￿￿ ￿￿￿￿￿￿￿ ￿￿￿￿￿ ￿￿ ￿￿￿￿￿￿￿￿￿￿￿￿￿ ￿￿￿￿￿￿￿￿￿￿￿￿￿ ￿￿ ￿￿￿￿￿￿￿ ￿￿ ￿￿￿￿￿￿￿￿￿ ￿￿￿￿￿ ￿￿￿￿￿ ￿￿￿￿￿￿￿ ￿￿￿￿￿. ￿￿￿￿￿￿￿￿￿￿, ￿￿￿￿￿￿￿￿ ￿￿ ￿￿￿￿￿￿￿￿￿ ￿￿￿￿’￿￿￿￿￿￿￿ ￿￿￿￿￿￿￿, ￿￿￿￿￿￿￿￿ ￿￿￿￿￿￿￿￿ ￿￿￿￿￿ ￿￿￿￿￿￿￿. ￿￿￿￿￿￿￿ ￿￿￿￿￿ ￿￿￿￿￿￿￿￿ ￿￿ ￿￿￿￿￿￿￿￿￿ ￿￿￿￿￿￿￿ ￿￿￿￿￿￿￿￿￿￿￿￿ ￿￿￿￿￿￿￿￿￿ α ￿￿￿￿￿ ￿￿￿￿￿￿￿. ￿￿￿￿￿￿￿￿￿￿ 1. Heath, Michael (2000). Scientific Computing. McGraw-Hill. ￿. 324. 2. ￿￿￿￿￿￿￿￿￿ ￿.￿. ￿￿￿￿￿￿￿￿ ￿￿￿￿￿￿￿￿ ￿ ￿￿ ￿￿￿￿￿￿￿￿￿ . ￿￿￿￿￿￿￿￿￿￿￿￿￿￿ ￿￿￿￿￿￿￿￿￿￿￿, 2013. ￿￿￿. 17. ￿. 93-106. 3. ￿￿￿￿￿ ￿.￿., ￿￿￿￿￿￿￿￿￿ ￿.￿. ￿￿￿￿ ￿￿￿￿￿￿￿ ￿￿￿￿￿￿￿￿ ￿￿￿￿￿￿￿￿￿￿ ￿￿￿￿￿￿ ￿￿￿￿￿￿￿￿ ￿￿￿￿￿￿￿￿￿￿￿ ￿￿￿￿￿￿￿￿￿￿￿￿. ￿￿￿￿￿￿￿ ￿￿￿￿￿￿￿￿ ￿￿￿￿￿￿￿￿￿￿￿ . ￿￿￿￿￿￿￿￿￿￿, 2015, ￿￿￿. 4. ￿.18-22. 165 4. O. Zalevska et al., "￿onstruction and study of the mathematical model for the 2 system using three -dimensional cellular automata," 2021 ieee 16th international conference on the experience of designing and application of cad systems (cadsm), 2021, pp. 49-52. 5. ￿￿￿￿￿￿￿￿￿, ￿. ￿. ￿￿￿￿￿￿￿￿￿￿￿￿ ￿￿￿￿￿￿￿￿￿￿￿￿￿￿ ￿￿￿￿￿￿￿ ￿￿￿￿￿. ￿￿￿￿’￿￿￿￿￿￿ ￿￿￿￿￿￿￿￿￿￿￿ ￿ ￿￿￿￿￿ ￿ ￿￿￿￿￿￿￿￿￿￿￿ ￿￿ ￿￿￿￿￿￿￿￿ ￿￿￿￿￿￿￿ ￿￿￿￿￿￿￿￿, ￿￿￿￿, 2014. ￿. 67-73. 6. ￿￿￿￿￿￿￿￿￿, ￿. ￿., ￿￿￿￿￿￿￿￿￿￿￿ ￿.￿. ￿￿￿￿￿￿ ￿￿￿￿￿￿ ￿￿￿￿￿￿￿￿￿ ￿￿￿￿￿￿￿￿￿￿￿￿￿￿￿ ￿￿￿￿￿￿￿ ￿￿￿￿￿ ￿￿ ￿￿￿￿￿￿￿￿￿￿￿￿ ￿￿￿￿￿￿￿￿￿￿￿ ￿￿￿￿￿￿- ￿￿. ￿￿￿￿￿￿￿ ￿￿￿￿￿￿￿￿ ￿￿￿￿￿￿￿￿￿￿￿. ￿￿￿￿￿￿￿￿￿￿, 2020. ￿￿￿.19.￿.138-146. 7. Sydorenko Yu. V ., Horodetskyi M.V . Modification of the algorithm for selecting a variable parameter of the Gaussian interpolation function. Control Systems and Computers, 2020, Issue 6 (290), pp. 21-28. ￿￿￿￿￿￿ ￿￿￿￿￿￿￿￿￿ ￿￿￿￿￿￿ ￿￿￿￿￿￿￿￿￿￿ ￿￿￿ ￿￿￿￿- ￿￿￿￿￿￿￿￿￿￿￿￿ ￿￿￿￿￿￿￿ ￿￿￿￿￿ ￿￿￿￿￿￿￿￿￿ ￿.￿., ￿￿￿￿￿￿￿￿￿ ￿.￿. ￿￿￿￿￿￿￿￿￿￿ ￿.￿., ￿￿￿￿￿￿￿￿￿ ￿.￿. ￿ ￿￿￿￿￿￿ ￿￿￿￿￿￿￿￿￿￿￿￿￿￿￿ ￿￿￿￿￿￿￿￿￿￿￿ ￿￿￿￿￿￿￿￿￿￿ ￿￿￿ ￿￿￿￿￿￿￿￿￿￿ ￿ ￿￿￿￿￿￿￿ ￿￿￿￿￿ ￿￿￿￿￿-￿￿￿￿￿￿￿￿￿￿￿￿. ￿￿￿￿￿￿￿￿￿￿ ￿￿￿￿￿￿￿￿￿￿￿ ￿￿￿￿￿￿￿￿￿￿ ￿￿ ￿￿￿￿ ￿￿￿￿￿￿￿￿￿￿￿￿￿￿￿ ￿￿￿￿￿￿￿ ￿￿￿￿￿￿￿￿￿￿￿￿ ￿￿￿￿￿￿￿￿￿￿￿￿. ￿￿￿￿￿￿￿￿￿￿ ￿￿￿￿￿￿￿ ￿￿￿￿￿￿ ￿￿￿￿￿￿￿￿￿￿￿￿ ￿￿￿￿￿￿￿￿￿ ￿￿￿￿￿￿￿￿￿￿￿￿ ￿￿￿￿￿ ￿￿￿￿￿￿ ￿￿￿￿￿- ￿￿￿￿￿￿￿ ￿ ￿￿￿￿￿￿￿￿￿￿￿￿ ￿￿￿￿￿￿￿￿￿￿ ￿ ￿￿￿￿￿￿￿￿￿￿￿￿￿ ￿￿￿￿￿￿ ￿￿￿￿￿￿￿￿. ￿ ￿￿￿￿￿￿￿￿￿ ￿￿￿￿￿ ￿￿￿￿￿￿￿￿ ￿￿￿￿￿￿￿￿￿￿￿￿￿ ￿￿￿￿￿￿￿￿￿￿￿￿ ￿￿￿￿￿ ￿￿￿￿￿￿￿￿ ￿￿￿￿￿￿￿￿￿, ￿￿ ￿￿￿￿￿￿￿￿￿￿￿￿ ￿￿￿￿￿￿￿￿￿￿￿￿ ￿￿￿￿￿￿￿ ￿￿ ￿￿￿￿￿￿￿￿￿￿￿￿, ￿￿￿￿￿￿￿￿￿￿ ￿￿￿￿￿, ￿￿￿￿￿￿￿￿ ￿￿￿, ￿￿￿￿￿￿￿￿￿ ￿￿￿￿￿￿￿￿ ￿ ￿￿￿￿￿￿￿ ￿ ￿￿￿￿￿ ￿￿￿￿￿￿￿￿￿ ￿￿￿￿￿￿ ￿ ￿￿ ￿￿￿￿￿￿￿￿ ￿ ￿￿￿￿￿￿￿ ￿ ￿￿￿￿￿￿￿￿￿. ￿￿￿￿￿￿￿ ￿￿￿￿￿￿￿￿￿￿￿￿ ￿ ￿￿￿￿￿￿￿￿￿￿￿ ￿￿￿￿￿￿￿￿￿￿￿￿￿￿￿￿ ￿ ￿￿￿￿￿ ￿￿￿￿￿￿￿￿￿￿￿ ￿￿￿￿￿￿￿￿ ￿￿￿￿￿ ￿￿￿￿￿￿￿￿￿￿￿￿ ￿￿￿￿￿. ￿￿￿￿￿￿￿ ￿￿￿ ￿￿￿￿￿￿￿ ￿￿￿￿￿￿￿￿￿ ￿￿￿￿￿￿￿￿. ￿ ￿￿￿￿￿￿￿￿￿￿￿ ￿￿￿￿￿￿￿￿￿￿￿￿￿￿￿ ￿￿￿ -￿￿￿￿￿￿ ￿￿￿￿￿￿￿￿ ￿￿￿￿￿￿￿￿ ￿￿￿￿￿￿￿￿￿￿, 2
\par
\[￿￿￿￿￿￿￿￿, 460 \times 800￿￿￿￿￿￿￿￿, ￿￿￿￿￿￿￿￿￿￿￿￿￿, ￿￿￿￿￿￿￿￿￿￿￿￿￿￿￿￿￿￿￿￿￿￿￿￿￿￿￿￿￿￿￿￿￿￿￿￿￿￿￿￿￿￿￿￿￿￿￿￿￿￿￿￿￿￿￿￿￿￿￿￿￿￿￿￿￿￿￿￿￿￿￿￿￿￿￿￿.\]
￿￿￿ ￿￿￿￿￿￿￿￿, ￿￿￿ ￿￿￿￿￿￿￿￿￿￿￿￿￿￿ ￿￿￿￿￿￿￿￿￿￿￿, ￿￿￿ ￿￿￿ ￿￿￿￿￿￿￿￿￿￿, ￿￿￿ ￿ ￿￿￿ ￿￿￿￿￿￿￿￿￿￿ ￿￿￿￿￿￿￿￿￿￿, ￿￿￿￿￿￿￿￿, ￿￿￿￿￿￿￿￿￿￿￿￿ ￿￿￿￿￿￿￿￿￿￿￿ ￿￿￿ ￿￿￿￿￿￿￿￿￿. ￿￿￿￿￿￿￿￿￿￿￿￿￿￿￿￿ ￿￿￿￿￿￿￿￿ ￿￿￿￿￿ ￿ ￿ ￿￿￿￿￿￿￿￿ ￿￿￿￿￿￿￿￿. ￿￿￿￿￿￿ ￿￿￿ ￿￿￿￿￿￿￿￿￿￿ ￿￿￿￿￿￿￿￿￿￿￿￿ ￿￿￿￿￿￿￿￿￿￿￿￿ ￿￿￿￿￿￿￿￿￿￿ ￿￿￿￿￿￿￿￿￿￿ ￿￿￿￿￿￿￿￿￿￿￿￿ ￿￿￿￿￿ ￿￿￿￿￿￿￿￿. ￿￿￿ ￿￿￿￿￿￿￿ ￿￿￿￿￿￿￿￿￿￿￿￿￿ ￿￿￿￿￿￿ ￿￿￿￿￿￿￿￿ ￿￿￿￿ ￿￿￿￿￿￿￿￿￿￿ ￿￿￿￿￿￿￿￿￿ ￿￿￿￿￿￿￿￿￿￿. ￿￿￿, ￿￿￿￿ ￿￿￿￿￿ ￿￿￿￿￿￿￿￿￿￿￿￿ ￿￿￿ ￿￿￿￿￿￿￿￿￿ ￿￿￿￿￿￿ ￿￿￿￿￿￿￿￿￿￿￿ ￿￿￿￿￿￿￿ ￿￿ ￿￿ 166 ￿￿￿￿￿￿￿￿￿￿ ￿￿￿￿￿￿￿￿￿, ￿￿ ￿￿￿ ￿￿￿￿￿￿￿￿￿ ￿￿￿￿￿￿￿￿￿￿￿￿￿￿￿￿￿ ￿￿￿￿￿￿ ￿￿ ￿￿￿￿￿￿ ￿￿￿￿￿￿￿￿￿￿￿￿￿ ￿￿￿ ￿￿￿￿￿￿￿￿￿￿￿￿. ￿￿￿￿￿ ￿￿￿￿￿￿￿￿￿￿￿￿ ￿￿￿￿￿￿￿￿ ￿￿￿￿￿￿ ￿￿￿￿￿￿￿￿￿￿￿ ￿￿￿￿￿￿￿￿￿￿￿￿￿ ￿￿￿￿￿￿ ￿￿￿￿￿￿￿￿ ￿￿ ￿￿￿￿￿ ￿￿ ￿￿￿￿￿￿￿ ￿￿￿￿￿. ￿ ￿￿￿￿￿￿ ￿￿￿￿￿￿￿￿￿￿ ￿￿￿￿￿￿￿￿￿￿￿￿￿￿￿ ￿￿￿ ￿￿￿￿￿￿￿ ￿ ￿￿￿￿￿￿￿ ￿￿￿￿ ￿￿￿￿￿ ￿￿￿￿￿￿￿ ￿￿￿￿￿￿, ￿￿￿￿￿￿￿￿ ￿￿￿￿￿￿￿￿￿￿￿￿ ￿￿￿￿￿ ￿￿￿￿￿￿￿ ￿￿￿ ￿￿￿￿￿￿￿ ￿￿￿￿￿￿￿￿￿￿￿￿ ￿￿￿￿￿￿ ￿￿￿￿￿￿￿￿￿￿￿￿￿￿￿￿ ￿￿￿￿￿￿￿ ￿￿￿￿￿. ￿￿￿￿￿￿￿￿ ￿￿￿￿￿: ￿￿￿￿￿￿￿￿￿￿￿￿, ￿￿￿￿￿￿￿￿￿￿￿￿￿￿￿￿ ￿￿￿￿￿￿￿ ￿￿￿￿￿￿, ￿￿￿￿￿￿￿￿￿￿￿ ￿￿￿￿￿￿￿￿￿￿￿￿, ￿￿￿￿￿￿￿ ￿￿￿￿￿. ANAL YSIS OF CALCULATION ERROR BEHA VIOR IN GAUSS INTERPOLATION OF RUNGE FUNCTION Sydorenko Iuliia, Zalevska Olga, Horodetskyi Mykola, Spirintsev Dmytro The paper considers the calculation error when ap plying Gaussian interpolation to the Runge function. The error is reduced by automatically selecting the variable coefficient. Examples of a computer program for interpolation of three Gaussian functions types are given and the results are compared with the errors of the Lagrange method. Nowadays, there are many examples of using interpolation, from the manufacture of mechanical parts in production, noise filtering, resampling, signal compression, etc. in the field of data processing and highlighting tren ds in the economy. Humanity is faced with the use of interpolation in the most unexpected areas of its daily lives. For example, the specifications of a phone's webcam have a resolution 460 × 800 pixels and it is also indicated that the camera can do image interpolation to achieve a higher resolution. This means that image conversion, both when zooming in and out, is done by interpolation. Interpolation also works with a color spectrum. Usually, for interpolation, classical literature suggests using the Lagrange method. The analysis of the use of the Lagrange method revealed some shortcomings. This method is used in the numerical estimation of the derivative function by its discrete values, but it is not always advisable to use it to process experimental data. The aim of the research is to analyze the disadvantages of using the Lagrange method on the example of the Runge function. The paper proposes to interpolate this function using three types of Gaussian function, the advantages of this approach for solving the classical Runge function interpolation problem are shown. Keywords: interpolation, Gaussian interpolation function, interpolation error, Runge phenomenon. 167 Referenses 1. Heath, Michael (2000). Scientific Computing. McGraw -Hill. ￿. 324. [in English] 2. Khovansky, AG (2013). Chebyshev polynomials and their inversions . Mathematical 3 education. 17, 93-106. [in Ukrainian] 3. Vanin, V .V ., Zalevskaya , O.V . (2015) Description of stable positions of dynamical systems by means of fractal approximation. Modern pro blems of modeling. Melitopol, 4, 18-22 [in Ukrainian] 4. O. Zalevska et al., (2021) "Construction and study of the mathematical model for the system using three -dimensional cellular automata," 2021 ieee 16th international conference on the experience of desi gning and application of cad systems (cadsm), 49-52. [in English] 5. Sidorenko, Yu. V . (2014) Parametric Gaussian interpolation function. Computer modeling in chemistry and technology and sustainable development systems KMHT, 67-73. [in Ukrainian] 6. Sidorenko, Yu.V ., Gorodetsky M.V . (2020) Analysis of the algorithm of the Gaussian interpolation function on elementary algebraic functions. Modern problems of modeling. 19, 138-146. [in Ukrainian] 7. Sydorenko Yu.V ., Horodetskyi M .V . (2020) Modification of the algorith m for selecting a variable parameter of the Gaussian interpolation function. Control Systems and Computers, 6 (290), 21-28. [in English] 4 3
\par
\end{document}
